\chapter{Scientific questions as access problems}
\label{foundations:chapter}

\section{Start with a quantity, a system, and an experiment}

A scientist rarely needs every microscopic parameter of a system. A useful
question is usually narrower: which direction does a current flow, how much
fuel does a process consume, does a particular error mechanism occur, or which
of several interventions changes a reaction pathway? The framework begins
with that quantity and asks whether the available experiments can determine
it. This order matters. A model can contain invisible parameters while the
quantity of interest remains identifiable.

Here is a small example that will recur in more detailed physical form.
Suppose two parallel channels contribute positive rates $a$ and $c$ to the
same observed transition. A detector identifies only the total rate
$k=a+c$. Replacing $(a,c)$ by $(a+h,c-h)$, while keeping both positive,
preserves every observation whose law depends only on $k$. The total rate
is identifiable; the individual rates are not. If a calibrated intervention
multiplies the first channel by a known $r\ne1$, a second observed total
$k_r=ra+c$ gives
\begin{equation}\label{foundations:twochannel}
 a=\frac{k-k_r}{1-r},\qquad c=\frac{k_r-rk}{1-r}.
\end{equation}
The intervention changes what can be identified. When $r$ is close to one,
the same formulas amplify measurement and calibration error. If $r$ is
unknown, the new observation need not remove the ambiguity. Thus an
experimental design requires both a model of intervention and a statement
of its accuracy.

The general object is an \emph{access specification}. It contains:
\begin{enumerate}
\item a candidate family $\Theta$, including nuisance parameters and the
      assumptions that make it a physically admissible family;
\item a target $g:\Theta\to\mathcal G$, with the units and required
      precision or decision loss stated;
\item available preparations, controls, measurements, retained records,
      and the rules governing their histories;
\item a calibration model and a resource account, including preparation,
      reset, sensing, control, and acquisition where these are used.
\end{enumerate}
This specification is a scientific input. Algebra can test its consequences;
it cannot establish from notation alone that an apparatus realizes it.
In particular, knowledge of the transition probabilities does not determine
the energy used by an unobserved subsystem.

Preparation belongs to access. A system that starts in an unknown state can
be observationally equivalent to a different system with another preparation.
A known reset, an independently calibrated stationary preparation, or a
recorded preparation outcome can change that conclusion. If preparation
depends on an unknown physical parameter, that dependence belongs in the
candidate model. It is not information available to the experimenter for
free.

\section{Transcript laws and targets}

Let $\mathcal A$ denote an access specification and $\Pi_{\mathcal A}$ its
permitted protocols. A protocol $\pi$ selects actions causally from retained
observations and any recorded or modeled randomization; it specifies when
to stop. Its complete accessible transcript has probability law
$P_\theta^\pi$ under candidate $\theta$. A transcript may contain action
choices, timestamps, detector outputs, preparation outcomes, and stopping
information. Omitting one of these changes the law being compared.
The same protocol is used under every candidate. A policy that secretly
knows the true candidate is not an experiment for distinguishing candidates.

\begin{definition}[Operational equivalence]\label{foundations:equivalence}
For a fixed access specification, write
\begin{equation}
 \theta\sim_{\mathcal A}\theta'
 \quad\Longleftrightarrow\quad
 P_\theta^\pi=P_{\theta'}^\pi
 \quad\text{for every }\pi\in\Pi_{\mathcal A}.
\end{equation}
The complete law profile is
$\Phi_{\mathcal A}(\theta)=(P_\theta^\pi)_{\pi\in\Pi_{\mathcal A}}$.
A target is \emph{structurally identifiable under $\mathcal A$} when
equality of these profiles implies equality of the target.
\end{definition}

This is identifiability from ideal experiment laws. It does not mean that
a finite noisy record determines a real parameter exactly. The distinction
is standard in statistical and dynamical-system identification
\cite{f_nemcova,f_thon}. It remains essential when the word ``observable''
is used with different conventions in different fields.

\begin{theorem}[Target factorization]\label{foundations:factorization}
A target $g$ is structurally identifiable under $\mathcal A$ if and only if
it is constant on each equivalence class of $\sim_{\mathcal A}$. Equivalently,
there is a unique map $\bar g$ on the attained profiles
$\Phi_{\mathcal A}(\Theta)$ such that
\begin{equation}
 g=\bar g\circ\Phi_{\mathcal A}.
\end{equation}
\end{theorem}
\begin{proof}
If a factorization exists, equal profiles have equal images under
$\bar g$ and hence equal targets. Conversely, define
$\bar g(\Phi_{\mathcal A}(\theta))=g(\theta)$. Constancy on each class
makes this definition independent of the representative. Every element
of the attained profile space has such a representative, proving uniqueness.
\end{proof}

The theorem is a set-theoretic quotient statement, not a claim that the
factor map is measurable, continuous, numerically stable, or efficiently
computable. Those properties need additional structure. Its practical use
is to separate an experimental obstruction from a difficult calculation.
An algorithm cannot reconstruct a target that takes two values on the
same attained profile without using further assumptions or further access.

For example, a hidden-state model may have several coordinate descriptions
related by a relabeling of latent states. Individual coordinates need not be
identifiable, while predictions and invariant target quantities are. The
quotient treats such representational ambiguity and physically different
but observationally equivalent mechanisms by the same mathematical rule.
Their scientific interpretations still differ.

\section{Processing, comparison, and the value of additional access}

For probability laws on a common measurable space define
\begin{equation}
 \operatorname{TV}(P,Q)=\sup_B|P(B)-Q(B)|.
\end{equation}
For a finite record alphabet this equals
$\frac12\sum_y|P(y)-Q(y)|$. A randomized processing rule is a Markov
kernel $K$: it maps a complete record to a reported statistic or decision
without knowing the true candidate.

\begin{proposition}[Data processing and access refinement]
\label{foundations:processing}
For any Markov kernel $K$,
\begin{equation}
 \operatorname{TV}(KP,KQ)\leq\operatorname{TV}(P,Q).
\end{equation}
Suppose every protocol $\pi$ available under $\mathcal A$ can be simulated
by some protocol $\rho$ under $\mathcal B$, followed by a
candidate-independent kernel $K$, so that
$P_\theta^\pi=K P_\theta^\rho$ for all $\theta$. Then
$\theta\sim_{\mathcal B}\theta'$ implies
$\theta\sim_{\mathcal A}\theta'$. Every target identifiable under
$\mathcal A$ is therefore identifiable under $\mathcal B$.
\end{proposition}
\begin{proof}
For any output event $B$, the function $z\mapsto K(z,B)$ takes values in
$[0,1]$. Integrating this function against $P-Q$ has absolute value at most
$\operatorname{TV}(P,Q)$, by the layer-cake representation as an integral
of indicator functions. Taking the supremum over $B$ proves contraction.
If the two $\mathcal B$ laws coincide, applying their common kernel gives
equal $\mathcal A$ laws for every $\pi$. Target identifiability then follows
from Theorem~\ref{foundations:factorization}.
\end{proof}

This is the familiar comparison-of-experiments principle
\cite{f_blackwell}, used here to compare experimental capabilities. A more
accurate detector, a new control, a retained controller record, or an allowed
reset can refine the equivalence classes. Conversely, discarding the action
history or averaging away a temporal correlation can merge classes.
The comparison needs a real simulation relation: two devices with
different resolution or resource requirements are not ordered merely
because one has more components.

For two fixed candidates and a declared resource budget $C$, let
\begin{equation}\label{foundations:frontier}
 D_{\mathcal A}(C;\theta_0,\theta_1)
 =\sup_{\pi\in\Pi_{\mathcal A}(C)}
   \operatorname{TV}(P_{\theta_0}^\pi,P_{\theta_1}^\pi).
\end{equation}
Here $\Pi_{\mathcal A}(C)$ must specify whether $C$ bounds elapsed time,
gross supplied energy, expected energy, operation count, or another
quantity. These constraints define different feasible sets. If a simulation
preserves the stated budget, Proposition~\ref{foundations:processing}
also orders the corresponding discrimination frontiers.

For known equally likely alternatives, the infimum average error over
protocols and tests is
\begin{equation}\label{foundations:bayes}
 \varepsilon^*(C)=\frac{1-D_{\mathcal A}(C;\theta_0,\theta_1)}{2}.
\end{equation}
Indeed, for a fixed protocol the likelihood comparison selects the larger
of the two probabilities or densities at each record, giving error
$(1-\operatorname{TV})/2$. Optimizing over protocols proves the expression.
Suprema and infima need not be attained. In particular, $D=1$ can describe
a limiting sequence of improving experiments without any single finite
experiment having zero error.

If two candidates with different targets are equivalent, every randomized
decision has the same distribution under both. A binary decision between
the two has average success at most $1/2$ under equal priors. For a real
target, every estimate $\widehat g$ satisfies
\begin{equation}\label{foundations:losslower}
 \max_{i=0,1}\mathbb E_{\theta_i}
       |\widehat g-g(\theta_i)|
 \geq \frac{|g(\theta_0)-g(\theta_1)|}{2}.
\end{equation}
The expectations use the same law for $\widehat g$; the triangle
inequality bounds the sum of the two risks below by the target separation.
Unlimited reprocessing of the same accessible evidence cannot remove this
obstruction. Additional physical access or a justified restriction of the
candidate family can.

\section{A common geometry for uncertainty and abstention}

Exact equality is only one source of ambiguity. Calibration bounds and
finite records also leave sets of compatible candidates. Let $z$ be reported
evidence and let a specified relation $\mathcal R\subseteq\Theta\times
\mathcal Z$ describe what that evidence permits. Define
\begin{equation}
 \Theta(z)=\{\theta:(\theta,z)\in\mathcal R\},\qquad
 G(z)=\{g(\theta):\theta\in\Theta(z)\}.
\end{equation}
For example, the relation may be exact equality of an observable profile,
a deterministic measurement enclosure, or a confidence construction with
a separately proved coverage guarantee. These interpretations should be
recorded, because they support different conclusions.

\begin{proposition}[Compatible-target radius]\label{foundations:radius}
Suppose $g$ is real valued, $\Theta(z)$ is nonempty, and estimates may take
any value in $\mathbb R$. If $G(z)$ is bounded, write
$a=\inf G(z)$ and $b=\sup G(z)$. Then
\begin{equation}
 \inf_{t\in\mathbb R}\sup_{\theta\in\Theta(z)}|t-g(\theta)|
 =\frac{b-a}{2},
\end{equation}
attained at $t=(a+b)/2$. If $G(z)$ is unbounded, every finite estimate has
infinite worst-case absolute error.
\end{proposition}
\begin{proof}
For bounded $G(z)$, the supremum loss is
$\max\{|t-a|,|t-b|\}$, whether or not the endpoints are attained. Its
minimum is the half-width, attained at the midpoint. If $G(z)$ is unbounded,
there are target values arbitrarily far from each fixed $t$.
\end{proof}

An empty compatible set is an inconsistency between evidence and the
declared model/calibration contract, rather than a uniquely identified
mechanism. A nonempty set on which $g$ is constant supports an exact target
within that contract even when several mechanisms survive. A wide target
set supports an interval or abstention. If $\Theta(Z)$ covers the true
candidate with probability at least $1-\delta$, the target set $G(Z)$
inherits that same coverage; the geometric proposition does not create the
coverage guarantee. For categorical targets or another loss function, the
appropriate decision problem replaces the real-valued midpoint.

This elementary geometry gives the common meaning of several later outputs:
an exact obstruction, a calibrated target interval, and a refusal to make an
unsupported decision. The linear intervention theorem specializes it to a
computable global worst-case radius. Sequential model checking supplies
coverage under adaptive actions and stopping. Neither step assumes that
surviving candidates are the whole of nature.

\section{When pairwise certificates give a finite scientific program}

\begin{theorem}[Finite target-separating suite]
\label{foundations:suite}
Let $\Theta=\{\theta_1,\ldots,\theta_K\}$ be a finite family. Suppose
the permitted elementary experiments have finite transcript alphabets and
any selected finite list can be executed on fresh preparations, independently
conditional on the same fixed candidate. Assume structural target
identifiability relative to these experiments. Then at most
$K(K-1)/2$ experiments suffice to form a suite whose joint law separates
every pair with different target values. If the individual models have
finite linear representations, witnesses can be synthesized with the finite
bound of Theorem~\ref{finite:span}.
\end{theorem}
\begin{proof}
For each unordered pair with distinct targets, identifiability supplies an
experiment with unequal laws. On a finite alphabet, some transcript event
has unequal probability. Choose one such experiment for each pair, dropping
duplicates. Fresh preparation gives a product joint law for the suite under
each fixed candidate. Its marginal for the chosen pair's experiment differs,
so the joint laws differ for that pair. There are at most $K(K-1)/2$ pairs.
The finite-representation assertion follows by applying
Theorem~\ref{finite:span} to each pair.
\end{proof}

The fresh-preparation assumption is consequential: concatenating probes on
one uncontrolled evolving specimen need not preserve their earlier
discriminating laws. The fixed-candidate qualification prevents a nuisance
parameter from being redrawn between trials merely to simplify a likelihood.
The theorem gives a finite separation program, not a finite-sample accuracy
claim or an efficient selection algorithm. Positive gaps can be small.
Section~\ref{workflow:suite} turns this existence result into a
calibration-aware experiment-cover workflow with explicit sample bounds.

\section{What is universal here}

The law-profile, comparison, and compatible-target constructions apply to
any well-defined family of experiments and targets. They are universal in
that mathematical sense. Finite-dimensional algorithms, Markovian dynamics,
quantum instruments, calibrated interventions, and energy envelopes impose
additional assumptions to make particular questions constructive.

The broad scientific contribution is an explicit chain from a question to
an admissible experiment and a checked conclusion. Its components connect
established statistical decision theory, system identification, control,
and physical modeling. This volume does not infer new conservation laws
from the shared formalism, and it does not equate an access limitation with
a general limit on computation. It provides a way to locate the unresolved
distinction, state the extra access that would remove it, and test whether
that access is practical.

The original boundary-and-arrow question is one instance: fix the allowed
local record, compare candidate dynamics with different time orientations,
and ask whether orientation is constant on the resulting equivalence
classes. A geometric boundary must be modeled as a physical restriction on
preparation, dynamics, coupling, or measurement. It cannot be silently
replaced by an unspecified loss of information. This operational entry point
preserves the physical question while allowing the same tools to address
reaction channels, sensor design, and quantum error hypotheses.
